\documentclass[12pt, a4paper]{article}

\usepackage[utf8]{inputenc}
\usepackage[T1]{fontenc}
\usepackage[brazil]{babel}
\usepackage{graphicx}
\usepackage{setspace}
\usepackage{geometry}
\usepackage[scaled]{helvet}
\usepackage{float}   
\usepackage[printonlyused]{acronym}


\renewcommand{\familydefault}{\sfdefault}

\geometry{
 a4paper,
 top=3cm,
 bottom=2cm,
 left=3cm,
 right=2cm
}

\begin{document}

% ------------------------------------------------------------------------
% CAPA
% ------------------------------------------------------------------------
\begin{titlepage}
    \begin{center}
        \textbf{COLÉGIO POLITÉCNICO BENTO QUIRINO}
        
        \vspace{3.5cm}
        
        JOÃO RAFAEL PENDEZZA MEDEIROS \\
        RAFAEL JOHANN GARRIDO \\
        JOÃO RICARDO FRAZÃO ROMANO ROCHA \\
        MATEUS BORGES NOGUEIRA \\
        JOÃO VITOR SOTERAS \\
        OTAVIO MATIAS MEDEIROS
        
        \vspace{4.5cm}
        
        \textbf{\large PENDEZZA PIZZA: \\ Documentação Técnica de Desenvolvimento}
        
        \vfill
        
        \textbf{CAMPINAS - SP} \\
        \textbf{2026}
    \end{center}
\end{titlepage}

% ------------------------------------------------------------------------
% FOLHA DE ROSTO
% ------------------------------------------------------------------------
\begin{titlepage}
    \begin{center}       
        \textbf{\large PENDEZZA PIZZA: \\ Documentação Técnica de Desenvolvimento}
        
        \vspace{2cm}
        
        \hspace{.45\textwidth}
        \begin{minipage}{.5\textwidth}
            \begin{spacing}{1.0}
                \small Trabalho de Conclusão de Curso apresentado ao Colégio Politécnico Bento Quirino como requisito parcial para obtenção do título de Técnico.
            \end{spacing}
        \end{minipage}
        
        \vfill
        
        CAMPINAS - SP \\
        2026
    \end{center}
\end{titlepage}

% Resumo
\newpage
\section*{Resumo}
Aqui você coloca o resumo do projeto... "O projeto Pendezza Pizza visa solucionar..."
\vspace{1cm}
\\ \textbf{Palavras-chave:} Pizzaria, Sistema, Automação.

% Resumo em inglês
\newpage
\section*{Abstract}
Aqui você coloca o resumo do projeto... "O projeto Pendezza Pizza visa solucionar..."
\vspace{1cm}
\\ \textbf{Palavras-chave:} Pizzaria, Sistema, Automação.

% Lista de Ilustrações
\newpage
\listoffigures

% Lista de Abreviaturas
\newpage
\section*{Lista de Abreviaturas e Siglas}

\begin{description}
    \item[API] \textit{Application Programming Interface} 
    \item[AWS] Amazon Web Services
    \item[EC2] \textit{Elastic Compute Cloud}
    \item[HTTP] \textit{HyperText Transfer Protocol}
    \item[JSON] \textit{JavaScript Object Notation}
    \item[MVC] \textit{Model-View-Controller}
    \item[OAUTH2] \textit{Open Authorization 2.0}
    \item[PEPI] \textit{Pendezza Pizza} 
    \item[S3] \textit{Simple Storage Service}
    \item[SES] \textit{Simple Email Service}
    \item[SQL] \textit{Structured Query Language}
\end{description}

% Sumário
\newpage
\tableofcontents

% ------------------------------------------------------------------------
% INÍCIO DO CONTEÚDO
% ------------------------------------------------------------------------

% Introdução
\newpage
\section{Introdução}

\subsection{Sobre nós}
Agora apresentaremos a tragetória da organização Pendezza Pizza bem como o que move-a como organização e os pilares que sustentam seu posicionamento no mercado gastrônomicos e seus valores

\subsubsection{Nascimento da Marca}
A marca nasceu em cima de uma piada sobre abrir uma pizzaria, onde foi unido o sobrenome de um dos integrantes com o produto inicialmente vendido. 

Idealizada e arquitetada pelo ainda estudande sensual do ensino médio João Rafael Pendezza Medeiros (Foto insana minha), Pendezza Pizza ou PePi foi desenvolvida com intuito de ser autosuficiente quando o assunto é gestão empresarial.

\subsubsection{Valores}
Os valores da PePi como organização é ajudar todo dono de restaurante a gerênciar seu negócio, desde relatórios de vendas até cardápios personalizados, a PePi entende a dor do cliente e por isso é a número um quando o assunto é uso de tecnologia alinhado a workshop de restaurantes.

Em poucas palavras posiciono a PePi como um workshop de restaurantes italianos, porém desenvolvendo essa ideia, eu afirmo que PePi é a melhor amiga do cozinheiro e do consumidor final, já que buscamos oferecer confiança, eficiência, atenção às dores dos consumidores, proatividade e sempre sujeita a mudanças

\subsubsection{Identidade Visual}
A identidade visual da PePi busca refletir a tradição gastrônomica italiana


\begin{figure}
    \centering
    \includegraphics[width=0.5\linewidth]{LogoPendezzaPizza.png}
    \caption{Logotipo oficial Pendezza Pizza}
    \label{fig:placeholder}
\end{figure}

\newpage 

\subsection{Objetivos}
Nesta seção, divide-se os objetivos que dão norte ao desenvolvimento do projeto, divididos entre a meta principal e as etapas necessárias para sua conclusão.

\subsubsection{Objetivos Gerais}
O objetivo geral deste trabalho consiste em desenvolver um software de gestão para pizzarias que integre controle de estoque, frente de caixa e inteligência de dados.

\subsubsection{Objetivos Específicos}
Para alcançar o objetivo geral, definiram-se as seguintes metas:
\begin{itemize}
    \item Implementar um sistema de autenticação via OAuth2;
    \item Desenvolver uma interface mobile intuitiva para o consumidor;
    \item Estruturar um banco de dados relacional para armazenamento de pedidos.
\end{itemize}

\newpage


% Engenharia do Projeto
\section{Engenharia do Projeto}
Esta é a seção principal técnica.

\subsection{Arquitetura}
Explique a arquitetura (MVC, Microsserviços, etc.).


\subsection{Modelagens do Sistema}
Abaixo, apresentamos o diagrama de fluxo do sistema, diagrama de classes, diagrama de .


\newpage

% Tecnologias Utilizadas
\section{Tecnologias Utilizadas}
Nesta seção, apresentamos todas as escolhas técnicas para viabilizar o projeto Pendezza Pizza, justificando cada escolha

\subsection{Back End}
Para o desenvolvimento da lógica de negócios e API do servidor, optou-se pela utilização do framework \textbf{Spring Boot} (linguagem Java).

A escolha do Spring Boot se deve pelo extenso tempo de mercado, comunidade ativa e inúmeros frameworks disponíveis, conforto e praticidade

Abaixo, detalhamos os componentes integrados à arquitetura do Back End:

\begin{description}
    \item[Java:] Linguagem de desenvolvimento do framework, utilizado em cada parte do backend, segue um modelo de simulação, focando em objetos e dados, o que o torna perfeito para manipular e controlar o acesso aos dados.
    
    \begin{figure}[H]
        \centering
        \fbox{AQUI ENTRA O LOGO}
        \caption{Logotipo oficial da linguagem Java}
        \label{fig:logo}
    \end{figure}

    \item[Spring Doc:] Principal biblioteca de documentação de API do projeto spring hoje em dia no formato OpenApi.
        
    \begin{figure}[H]
        \centering
        \fbox{AQUI ENTRA O LOGO}
        \caption{Logotipo oficial da biblioteca Spring Doc OpenApi}
        \label{fig:logo}
    \end{figure}
    
    \item[Spring Security e OAuth2:] Utilizados para a camada de segurança. O protocolo OAuth2 garante que a autenticação e autorização dos usuários sejam feitas via tokens de acesso padronizados, permitindo um acesso controlado e seguro à aplicação.
    
    \item[Infraestrutura {AWS}:] O sistema está hospedado em ambiente de nuvem para garantir escalabilidade.
    \begin{itemize}
        \item \textbf{EC2:} Servidor virtual onde a aplicação Spring Boot está rodando.
        \item \textbf{S3:} Utilizado para armazenar as imagens das pizzas e do cardápio, retirando a carga do banco de dados.
        \item \textbf{SES:} Serviço utilizado para disparo automático de e-mails transacionais (confirmação de cadastro e pedidos).
    \end{itemize}
    
    \item[Monitoramento e Logs:] Implementou-se um sistema de logs para rastreabilidade de erros e auditoria das operações realizadas pelos usuários.
\end{description}

\subsection{Banco de Dados}

\subsection{Front End}

\subsection{QA/Testes}

\subsection{Inteligência Artificial}

\subsection{Softwares Usados}

\newpage

% Pesquisa de Mercado
\section{Pesquisa de mercado}
Resultados e Testes.

\newpage

% Bibliografia
\section{Bibliografia}

\end{document}